%!TEX root = ../thesis.tex
\chapter{系统的仿真与测试}
上文介绍了多机器人协同搜索的关键模块与技术方案，本章将介绍对此系统的仿真与测试效果。将综合验证系统的控制、搜索以及路线规划等能力。
\section{仿真环境}
\subsection{主要运行环境}
仿真环境的基本情况包括硬件配置和软件环境（表6-1），由于gazebo仿真环境涉及较大量的图像处理，对硬件性能和软件版本有一定需求，因此仿真环境的配置是值得关注的问题。

\begin{table}[ht]
    \caption{主要运行环境配置}
    \centering
    \resizebox{\textwidth}{!}{
    \begin{tabular}{lc} % 限制第二列宽度
        \toprule
        \textbf{项目} & \textbf{参数} \\
        \midrule
        内存 & 16GB \\
        处理器 & 1th Gen Intel® Core™ i7-11800H @ 2.30GHz × 16  \\
        显卡 & NVIDIA GeForce RTX 3060 Laptop GPU \\
        系统 & Ubuntu 20.04.6 LTS \\
        Gazebo版本 & version 11.15.1\\
        ROS 版本 & ros2-galactic\\
        \bottomrule
    \end{tabular}
    }
    \label{tab:laserscan}
\end{table}
在良好的配置下gazebo可以提高物理仿真参数以提高仿真的效果，使其更加接近实用条件，否则则有可能出现意料之外的错误效果。

\subsection{主要仿真参数}
对系统进行仿真时需要定义机器人参数，环境参数等，以求模拟接近真实环境。首先搭建了机器人进行搜索的地图，模拟现实环境下的房屋环境（图6-1），包括6个房间。其中有书架、桌子、垃圾桶、书架、啤酒罐等物品，所有材质都具有对应的纹理，例如木纹、大理石、塑料、红砖墙等，其中房屋结构有门、窗、走廊等。该房屋环境中共放置了5个啤酒罐模型，作为机器人搜索的目标。
\begin{figure}[ht] % use float package if you want it here
    \centering
    \includegraphics[width=0.8\textwidth]{房间仿真.png}
    \caption{仿真房屋环境搭建}
    \label{fig:ipe}
\end{figure}

在仿真环境中需要设置物理引擎参数（表6-2），物理引擎参数会决定仿真的效果和性能。本项目在硬件限制条件下尽可能调整real\_time\_update\_rate和max\_step\_size以提高时间准确性。

\begin{table}[ht]
    \centering
    \renewcommand{\arraystretch}{0.9}
    \caption{Gazebo物理引擎参数设置}
    \begin{tabular}{l l l}
        \toprule
        \textbf{参数} & \textbf{含义} & \textbf{值} \\
        \midrule
        real\_time\_update\_rate & 仿真每秒更新次数 & 1000.0 \\
        max\_step\_size & 每步最大仿真时间步长（秒） & 0.001 \\
        real\_time\_factor & 仿真时间与真实时间的比例 & 1 \\
        type & ODE 求解器类型 & quick \\
        iters & 求解器最大迭代次数 & 150 \\
        sor & 超松弛参数（加速收敛） & 1.4 \\
        use\_dynamic\_moi\_rescaling & 是否启用动态惯性矩缩放 & 1 \\
        cfm & 约束柔性参数（降低约束振荡） & 0.00001 \\
        erp & 误差恢复参数（影响刚体恢复速度） & 0.2 \\
        contact\_max\_correcting\_vel & 最大接触点修正速度 & 2000.0 \\
        contact\_surface\_layer & 接触表面层厚度（影响碰撞检测精度） & 0.01 \\
        \bottomrule
    \end{tabular}

    \label{tab:ode_parameters}
\end{table}

\section{在仿真环境下运行系统}
完成环境搭建后在仿真房屋的两个不同房间里分别放入turtlebot\_waffle机器人（图6-2），其中蓝色部分指示机器人激光雷达范围，白色线条指示机器人摄像头视角范围。

\begin{figure}[H] % use float package if you want it here
    \centering
    \includegraphics[width=0.7\textwidth]{放机器人1.png}
    \caption{在仿真房屋中放入机器人}
    \label{fig:ipe}
\end{figure}

运行地图融合模块，可以在rviz中查看当前根据两个机器人激光雷达获取得到的融合地图（图6-3）。
\begin{figure}[H] % use float package if you want it here
    \centering
    \includegraphics[width=0.7\textwidth]{地图融合.png}
    \caption{地图融合模块}
    \label{fig:ipe}
\end{figure}

运行yolo目标检测，可以看到yolov5开始对目标进行检测（图6-4）。
\begin{figure}[H] % use float package if you want it here
    \centering
    \includegraphics[width=0.7\textwidth]{yolo运行.png}
    \caption{目标检测模块运行信息}
    \label{fig:ipe}
\end{figure}

运行目标搜索与定位模块，rviz中将会显示三个窗口，分别展示两个机器人的摄像头画面和标记了主要信息的地图（图6-5）。其中机器人的摄像头画面经过目标检测模块的标记，带有目标检测框信息（图6-6）。地图展示窗口中的标记包括目标坐标标记（图6-7）和规划路径标记，两个机器人的路径和搜索目标坐标分别标记为红色和绿色（图6-8）以便区分。其中较大的亮黄色点即指示数据处理后的目标位置（图6-9）。
\begin{figure}[H] % use float package if you want it here
    \centering
    \includegraphics[width=0.8\textwidth]{目标搜索.png}
    \caption{目标搜索与定位显示界面}
    \label{fig:ipe}
\end{figure}

\begin{figure}[H]
    \centering
    \begin{minipage}{0.4\textwidth} 
        \centering
        \includegraphics[width=\textwidth]{目标检测框.png}
        \caption{带有目标检测框的图像}
    \end{minipage}
    \begin{minipage}{0.4\textwidth} 
        \centering
        \includegraphics[width=\textwidth]{目标坐标点.png}
        \caption{原始目标坐标点}
    \end{minipage}
    \begin{minipage}{0.4\textwidth} 
        \centering
        \includegraphics[width=\textwidth]{规划路线修.png}
        \caption{红绿色规划路线}
    \end{minipage}
    \begin{minipage}{0.4\textwidth} 
        \centering
        \includegraphics[width=\textwidth]{数据处理后坐标.png}
        \caption{数据处理后坐标}
    \end{minipage}

    \label{fig:minipage_comparison}
\end{figure}

运行路径规划与控制模块，可以看见机器人开始请求路径并且模块响应，然后机器人开始运动。

\begin{figure}[H] % use float package if you want it here
    \centering
    \includegraphics[width=0.8\textwidth]{路径规划响应修.png}
    \caption{路径规划模块响应请求}
    \label{fig:ip}
\end{figure}

\section{系统测试}
在仿真环境下运行系统后，机器人逐渐完成对该环境的搜索（图6-11），在所有边界被探索完成后，路径规划与运动控制模块会发出探索完毕的信息。最终得到一张该环境的完整地图，并带有目标的坐标（图6-11）。

\begin{figure}[H] % use float package if you want it here
    \centering
    \includegraphics[width=0.8\textwidth]{探索完成.png}
    \caption{路径规划模块完成探索}
    \label{fig:ip}
\end{figure}

\begin{figure}[H] % use float package if you want it here
    \centering
    \includegraphics[width=0.8\textwidth]{系统运行结果.png}
    \caption{运行完毕得到目标地图}
    \label{fig:ip}
\end{figure}

观察该地图中的黄色目标标记，与仿真环境中的目标分布相对比，可见搜索结果与实际情况基本对应，找出了全部啤酒罐目标，数据处理后的目标坐标与真实位置平均偏差不到0.7米。但是也可以发现出现了一个多余的误报目标坐标，推测是由于环境中桌腿阻挡导致的激光雷达数据异常。除此之外，系统建立的地图完整清晰，机器人协同运行流畅无冲突。测试时，由15时36分53机器人开始请求路径到15时44分47秒完成搜索，完成探索共用时7分54秒，完成了对约90平方米的房屋环境搜索。至此可以认为，系统的全部功能已经实现。
